\begin{center}[13.]\par\medskip\textit{Fundamental Law Governing All Interactions of Galvanic Currents.}\par\smallskip(Found in July 1835.)\end{center}
Two elements of electricity that move relative to each other attract or repel each other differently than when they are at rest relative to each other.
\SourceLeaf{623}\[
\begin{gathered}
\begin{aligned}
e,x,y,z&\qquad\text{Element and its coordinates}\\
e',x',y',z'&
\end{aligned}\\[4pt]
(x'-x)^2+(y'-y)^2+(z'-z)^2=rr
\end{gathered}
\]
Interaction (repulsion)
\[
=\frac{ee'}{rr}\left\{1+k\left(\begin{gathered}
\left(\frac{\mathrm d(x'-x)}{\mathrm dt}\right)^2+\left(\frac{\mathrm d(y'-y)}{\mathrm dt}\right)^2\\
+\left(\frac{\mathrm d(z'-z)}{\mathrm dt}\right)^2-\tfrac{3}{2}\left(\frac{\mathrm dr}{\mathrm dt}\right)^2
\end{gathered}\right)\right\}
\]
where $\BareRoot{\frac{1}{k}}$ stands for a definite velocity.

\begin{center}[14.]\end{center}
In \textit{another} form, the \textit{fundamental law} reads as follows.

Let $P$ and $P'$ be points in two currents, with $x,y,z$ and $x',y',z'$ their coordinates and $r$ their distance apart. Let $\mathrm ds,\mathrm ds'$ be two current elements starting at those points.
\[
\begin{array}{rllll}
u&\text{Angle}&\text{between}&\mathrm ds\text{ and }\mathrm ds'\\
q&\text{-}&\text{-}&\mathrm ds\text{ and }PP'\\
q'&\text{-}&\text{-}&\mathrm ds'\text{ and }P'P
\end{array}
\]

If the factors that depend on the intensities of the currents are omitted, the elements $\mathrm ds,\mathrm ds'$ exert a mutual attraction that can be measured by
\[
\frac{\mathrm ds\,.\,\mathrm ds'\,.\,(\cos u+\tfrac{3}{2}\cos q\,.\,\cos q')}{rr}
\]

If we set $\mathrm dx=\mathrm dy=0$, this force is
\[
=\frac{\mathrm ds\,.\,\mathrm dz'}{rr}-\tfrac{3}{2}\,.\,\frac{z'-z}{r^3}\mathrm ds\,.\,\mathrm ds'\,.\,\left(\frac{\mathrm dr}{\mathrm ds'}\right)
\]
or, in terms of the component forces acting on $\mathrm ds$ parallel to the coordinate axes,
\[
\begin{gathered}
\mathrm ds\,.\,\left\{\frac{(x'-x)\mathrm dz'}{r^3}-\frac{3(x'-x)(z'-z)\mathrm d'r}{2r^4}\right\}\\[4pt]
\mathrm ds\,.\,\left\{\frac{(y'-y)\mathrm dz'}{r^3}-\frac{3(y'-y)(z'-z)\mathrm d'r}{2r^4}\right\}\\[4pt]
\mathrm ds\,.\,\left\{\frac{(z'-z)\mathrm dz'}{r^3}-\frac{3(z'-z)(z'-z)\mathrm d'r}{2r^4}\right\}
\end{gathered}
\]
\SourceLeaf{624}or, if one adds the forces
\[
\begin{gathered}
-\mathrm ds\,.\,\mathrm d'\frac{(x'-x)(z'-z)}{2r^3}\\[4pt]
-\mathrm ds\,.\,\mathrm d'\frac{(y'-y)(z'-z)}{2r^3}\\[4pt]
-\mathrm ds\,.\,\mathrm d'\frac{(z'-z)(z'-z)}{2r^3}
\end{gathered}
\]
which is permissible insofar as $\mathrm ds'$ is an element of a closed current, the components become
\[
\begin{gathered}
\frac{\mathrm ds}{2r^3}((x'-x)\mathrm dz'-(z'-z)\mathrm dx')\\[4pt]
\frac{\mathrm ds}{2r^3}((y'-y)\mathrm dz'-(z'-z)\mathrm dy')\\[4pt]
\frac{\mathrm ds}{2r^3}((z'-z)\mathrm dz'-(z'-z)\mathrm dz')
\end{gathered}
\]
The last is clearly $=0$. The following forces are clearly equivalent to these as far as the action on $\mathrm ds$ is concerned.
\[
\begin{aligned}
\text{1. In the direction }PP'\; .\;.\;.\;.\;.\quad&\frac{\mathrm ds\,.\,\mathrm dz'}{2rr}=\frac{\mathrm ds\,.\,\mathrm ds'}{2rr}\cos u\\[4pt]
\text{2. in the direction parallel to }\mathrm ds'\; .\;.\;.\;.\;.\quad&-\frac{\mathrm ds\,.\,\mathrm ds'}{2rr}\,.\,\cos q
\end{aligned}
\]
to which one may also add
\[
\text{3. in the direction }\mathrm ds\; .\;.\;.\;.\;.\;.\quad+\frac{\mathrm ds\,.\,\mathrm ds'}{2rr}\,.\,\cos q'
\]
while three forces exactly opposite to these will then act on $\mathrm ds'$.

Let $x,y,z$ be the coordinates of the acting current element and $0,0,0$ those of the element acted on. If the direction and strength of the first and second elements, resolved along the coordinate axes, are $\xi',\eta',\zeta';\ \xi,\eta,\zeta$, then according to \textsc{Ampère} the total force of attraction is
\[
\frac{\xi\xi'+\eta\eta'+\zeta\zeta'}{rr}-\tfrac{3}{2}\frac{(x\xi+y\eta+z\zeta)(x\xi'+y\eta'+z\zeta')}{r^4}
\]
so one component of the force is
\[
\begin{aligned}
&\frac{x}{r^3}(\xi'\mathrm dx+\eta'\mathrm dy+\zeta'\mathrm dz)-\frac{3x}{2r^4}(x\xi'+y\eta'+z\zeta')\mathrm dr\\[5pt]
&=\tfrac{1}{2}\xi'\mathrm d\frac{xx}{r^3}+\tfrac{1}{2}\eta'\mathrm d\frac{xy}{r^3}+\tfrac{1}{2}\zeta'\mathrm d\frac{xz}{r^3}\\
&\qquad+\frac{\eta'}{2r^3}(x\mathrm dy-y\mathrm dx)+\frac{\zeta'}{2r^3}(x\mathrm dz-z\mathrm dx)
\end{aligned}
\]
\SourceLeaf{625}or, since one can omit exact differentials,
\[
\frac{\eta'(x\eta-y\xi)+\zeta'(x\zeta-z\xi)}{2r^3}=\frac{x(\eta\eta'+\zeta\zeta')-y\xi\eta'-z\xi\zeta'}{2r^3}
\]
Here one may also add or omit
\[
\frac{x\xi\xi'+y\eta\xi'+z\zeta\xi'}{2r^3}
\]
which makes the formula symmetric in the two elements. If we choose the latter, we obtain the force
\[
\frac{x(-\xi\xi'+\eta\eta'+\zeta\zeta')-y(\xi\eta'+\eta\xi')-z(\xi\zeta'+\zeta\xi')}{2r^3}
\]
This is explained by a force that depends on the \textit{relative} motions $\alpha,\beta,\gamma$
\[
=-\frac{x(\alpha\alpha+\beta\beta+\gamma\gamma)-2\alpha(x\alpha+y\beta+z\gamma)}{2r^3}
\]

Here we can make do with a single additional force. It acts along $\mathrm{eb}$ with strength $\frac{vv}{kk}ee$ if $\mathrm{ea}$ is the direction of the relative motion $v$.
\[
\mathrm{aeb}=180^0-\mathrm{aeE},\qquad\mathrm{eE}\text{ positive}
\]
\begin{center}\begin{tikzpicture}[x=.115mm,y=.115mm,line width=.65pt,>=stealth]
  \coordinate (E) at (0,0);
  \coordinate (e) at (551,0);
  \coordinate (back) at (434,-77);
  \coordinate (a) at (780,138);
  \coordinate (b) at (685,222);
  \draw (E)--(e);
  \draw[->] (back)--(e)--(a);
  \draw[dash pattern=on 3pt off 2.4pt,->] (e)--(b);
  \filldraw[fill=white] (E) circle[radius=2.1pt];
  \filldraw[fill=white] (e) circle[radius=2.1pt];
  \node[left=3pt,inner sep=0pt] at (E) {\bfseries E};
  \node[anchor=north west,yshift=-3pt,inner sep=0pt] at (e) {\bfseries e};
  \node[above=2pt,inner sep=0pt] at (a) {\bfseries a};
  \node[above=2pt,inner sep=0pt] at (b) {\bfseries b};
\end{tikzpicture}
\end{center}

The law can be stated even more simply and in full generality as follows.
\begin{center}\begin{tikzpicture}[x=.115mm,y=.115mm,line width=.65pt,>=stealth]
  \coordinate (E) at (0,0);
  \coordinate (A) at (428,139);
  \coordinate (B) at (796,17);
  \coordinate (C) at (433,-124);
  \draw (E)--(B);
  \draw[->] (E)--(A);
  \draw[dash pattern=on 3pt off 2.4pt,->] (E)--(C);
  \filldraw[fill=white] (E) circle[radius=2.1pt];
  \node[anchor=south east,xshift=-3pt,inner sep=0pt] at (E) {\bfseries E};
  \node[above=2pt,inner sep=0pt] at (A) {\bfseries A};
  \node[right=4pt,inner sep=0pt] at (B) {\bfseries B};
  \node[below=2pt,inner sep=0pt] at (C) {\bfseries C};
\end{tikzpicture}
\end{center}

If an electric element $\mathrm E$ is acted on by another with a force in the direction $\mathrm{EA}$ and of strength $p$ when the two are at rest relative to one another, then in the case of relative motion whose direction%
\SourceLeaf{626}%
 with respect to $\mathrm E$ is the straight line $\mathrm{EB}$ and whose velocity is $=v$, a second force is added to the first, with
\[
\text{Strength}\quad=\frac{pvv}{kk},\qquad\text{Direction}\quad=\mathrm{EC}
\]
where $\mathrm{EA,EB,EC}$ lie in a single plane and $\mathrm{EB}$ is midway between $\mathrm{EA}$ and $\mathrm{EC}$.

